\paragraph{Sequence-level distillation of physical targets.}
Following the sequence-level distillation principle of Kim and Rush
(EMNLP 2016, doi:10.18653/v1/D16-1139), we drew a complete pictured
expression before computing its successive physical losses. For training
context $x$, the teacher assigned mass proportional to
$\exp(u(x)^\top v(y)/0.02)$ to each of 3,805 pictured training expressions
$y$, using normalized context and answer embeddings. At each epoch the
target was drawn from this distribution with probability 0.75 and was
the source expression otherwise. Both fits started from the same
checkpoint, froze the reader and ink renderer, and used one target per
passage per epoch. We verified identical cached and direct logits at
the same batch shape and verified equality of all frozen tensors after
training. With a batch size of 192 and AdamW learning rate 0.0006, the
two fits took 43.0 and 36.6 seconds. Combined fitting, assessment and
teacher preparation took 98.1 seconds on one RTX 4090~D with 335~MiB
peak allocated device memory. Earlier reader and renderer training are
additional costs. Sampling changed 1,212,654 of 1,714,536 targets.
Neither the option images nor the scoring head is needed to write.

For confirmation we took four unused source-group hash buckets (first
byte 32--35), excluding examples with complete pictured windows shared
with training, development or earlier confirmation examples. The
collection draws on the official training source, in the same genre
as development. The official test split was left unopened. Neighbour
ablations set the indicated slots to the all-white padding image and
preserved positions. Intervals resample complete source groups 2,000
times. Exact scoring requires both physical equality and the correct stop.
